\autoref{fig:approach-clustering} uses a schematic resource-feature cloud to illustrate the assignment rules. The cloud has variations around the Frontline and Specialist resolution patterns from the running example, along with an isolated observation.

\begin{figure}[htbp]
  \centering
  \subfloat[$k$-means\label{fig:approach-clustering-kmeans}]{%
    \FigurePlaceholder[0.28]{TikZ sketch of the shared points, fitted centroids, and nearest-centroid assignments.}}
  \hfill
  \subfloat[Self-organizing map\label{fig:approach-clustering-som}]{%
    \FigurePlaceholder[0.28]{TikZ sketch of the shared points linked to their winning cells on a grid with neighboring prototypes.}}
  \hfill
  \subfloat[DBSCAN\label{fig:approach-clustering-dbscan}]{%
    \FigurePlaceholder[0.28]{TikZ sketch of core neighborhoods and border points. Mark the isolated observation as noise.}}
  \caption{State assignment on a shared schematic feature cloud. Each point represents a resource feature window. The panels show the nearest centroid for $k$-means, the winning cell for a self-organizing map, and density-connected groups and noise for DBSCAN.}
  \label{fig:approach-clustering}
\end{figure}
